\chapter{动态抓取任务与专家控制器}\label{chap:task-expert}

\section{任务设置}

本文使用Genesis仿真平台构建动态立方体抓取任务\citep{genesis2024}。场景包含固定基座的Franka Research 3七自由度机械臂、Franka Hand二指夹爪\citep{franka2026research3,franka2026hand}和一个边长$0.05\,\mathrm{m}$的立方体。每个回合开始前，立方体在预设空间区域内随机生成，并以给定的线速度和角速度开始运动。机械臂需依次完成五个阶段——接近、对齐、闭合抓取、保持和搬运——最终将立方体送入目标区域。仿真步长为$\Delta t=0.01\,\mathrm{s}$；每个回合最多运行2400个仿真步。

任务目标是将立方体运送至半径$0.10\,\mathrm{m}$的目标区域内；本文以立方体中心到目标中心距离不超过$0.05\,\mathrm{m}$作为成功判据，达到该条件时记为成功回合。本文以闭环成功率为主指标，即成功回合数除以总评估回合数。主实验中，每个策略与场景组合均评估750个闭环回合。

\begin{figure}[!htbp]
    \centering
    \includegraphics[width=0.90\textwidth]{fig_task_setup}
    \caption{Genesis仿真平台中的FR3与Franka Hand动态立方体抓取任务。图中展示立方体生成区域和目标区域。}
    \label{fig:task-setup}
\end{figure}

\section{实验场景}\label{sec:experiment-scenes}

本文设置两个实验场景：定轴自旋场景（立方体姿态固定，仅绕$z$轴自旋）和翻滚场景（立方体初始姿态均匀随机，绕空间轴旋转）。两个场景的生成区域、线速度范围、机械臂配置和成功判据完全一致，关键区别在于抓取面是否预先确定：定轴自旋场景中抓取面固定且已知；翻滚场景中抓取面需在执行过程中动态确定。

两个场景的目的不同。定轴自旋场景是验证基线：由于抓取面固定且已知，该场景用于确认策略具备基本的轨迹跟踪和时序判断能力，为分析翻滚场景提供参照。翻滚场景是主实验场景：抓取面选择、姿态对齐和闭合时机判断在此才构成真实挑战——这正是本文研究问题的核心。

\subsection{定轴自旋场景}

立方体初始姿态固定（与世界坐标系原点姿态一致），绕世界坐标系$z$轴以$30$--$60\,\mathrm{deg/s}$的角速度自旋。线速度在$[-0.10,-0.10,-0.02]$到$[0.10,0.10,0.02]\,\mathrm{m/s}$的三维盒形范围内均匀采样，生成位置范围为$[0.20,-0.45,0.20]$到$[0.48,0.45,0.40]\,\mathrm{m}$。该场景中，控制器可预先确定抓取面，任务难点主要是轨迹跟踪和闭合时机判断。

\subsection{翻滚场景}

立方体初始姿态均匀随机（四元数在$\mathbb{S}^3$上均匀采样），绕空间轴以$15\,\mathrm{deg/s}$的角速度旋转。本文中”翻滚”指立方体绕空间轴旋转（初始姿态随机），并非物体在地面上接触滚动。线速度在$[-0.075,-0.075,-0.02]$到$[0.075,0.075,0.02]\,\mathrm{m/s}$的三维盒形范围内均匀采样，生成位置范围与定轴自旋场景相同。翻滚场景的核心挑战不在于自旋速度——其角速度$15\,\mathrm{deg/s}$甚至低于定轴场景的$30$--$60\,\mathrm{deg/s}$。真正的难度来自随机初始姿态和空间轴旋转的叠加：抓取面选择、姿态对齐和闭合时机判断这三个挑战相互耦合，是翻滚场景的核心难点。

综上，定轴自旋与翻滚场景构成由简至繁的两层测试环境。两个场景各采集5000条成功专家示教，共享相同的数据预处理流程（筛选、采样步幅、验证集比例等详见\S\ref{sec:expert-demos}）。

图~\ref{fig:ch2-scene-motion-snapshots}用同一相机视角展示两类场景中的立方体运动。为清晰展示立方体运动趋势，图中固定了起始位置并采用各场景允许的最大线速度与最大角速度。定轴自旋场景从固定初始姿态出发，抓取面不随回合变化；翻滚场景从随机初始姿态出发并绕空间轴旋转，姿态和可达抓取面随时间共同变化。

\begin{figure}[!htbp]
    \centering
    \captionsetup[subfigure]{font=footnotesize,skip=2pt}
    \begin{subfigure}[t]{0.31\textwidth}
        \includegraphics[width=\linewidth]{fig_ch2_scene_motion_0_0}
        \caption{定轴自旋，$t=0.0\,\mathrm{s}$}
    \end{subfigure}\hfill
    \begin{subfigure}[t]{0.31\textwidth}
        \includegraphics[width=\linewidth]{fig_ch2_scene_motion_0_1}
        \caption{定轴自旋，$t=0.8\,\mathrm{s}$}
    \end{subfigure}\hfill
    \begin{subfigure}[t]{0.31\textwidth}
        \includegraphics[width=\linewidth]{fig_ch2_scene_motion_0_2}
        \caption{定轴自旋，$t=1.6\,\mathrm{s}$}
    \end{subfigure}

    \vspace{0.45em}
    \begin{subfigure}[t]{0.31\textwidth}
        \includegraphics[width=\linewidth]{fig_ch2_scene_motion_1_0}
        \caption{翻滚，$t=0.0\,\mathrm{s}$}
    \end{subfigure}\hfill
    \begin{subfigure}[t]{0.31\textwidth}
        \includegraphics[width=\linewidth]{fig_ch2_scene_motion_1_1}
        \caption{翻滚，$t=0.8\,\mathrm{s}$}
    \end{subfigure}\hfill
    \begin{subfigure}[t]{0.31\textwidth}
        \includegraphics[width=\linewidth]{fig_ch2_scene_motion_1_2}
        \caption{翻滚，$t=1.6\,\mathrm{s}$}
    \end{subfigure}
    \caption{定轴自旋与翻滚场景的运动快照。上排为定轴自旋场景，下排为翻滚场景；三列对应同一回合中的三个时间点。}
    \label{fig:ch2-scene-motion-snapshots}
\end{figure}

\section{专家策略与示教采集}\label{sec:expert-demos}

专家策略由有限状态机（FSM）和MPC组成，包含五个阶段：(1)~接近阶段——根据预测的未来立方体姿态选择可达抓取面并移动至预抓取位姿；(2)~对齐阶段——持续跟踪短时抓取参考，待相对位姿误差满足闭合条件后转入闭合；(3)~闭合阶段——按固定时长闭合夹爪；(4)~保持阶段——维持抓取稳定；(5)~搬运阶段——将已抓取立方体运送至目标区域。FSM负责阶段切换、抓取时机和几何参考构造，MPC负责带约束的任务空间跟踪。

该专家既定义了示教分布，也作为第四章中模型控制方法的对比基线。本节从任务执行角度说明专家控制器的工作机制。专家控制器在每个回合开始时，根据当前末端状态和立方体真实状态生成一次抓取计划。计划使用$4.0\,\mathrm{s}$预测窗口，将时间分配给接近、对齐和闭合三个抓取相关阶段：对齐和闭合的预设时长分别为$0.80\,\mathrm{s}$和$0.30\,\mathrm{s}$；接近阶段时长取$0.50\,\mathrm{s}$与可用时间的一半中的较大值。

抓取面从立方体六个面中选择：控制器先用当前线速度和角速度外推接近结束时的立方体位姿，再对每个候选面计算可达性和朝向评分，排除不可达面后，综合考虑机械臂工作空间、面法向朝向和目标速度方向选定。选定抓取面后，专家生成两个几何参考：抓取面法向外侧$0.08\,\mathrm{m}$处的预抓取位姿，以及更靠近物体中心（沿法向进入物体内部$0.035\,\mathrm{m}$）的抓取位姿。接近阶段到达预抓取位姿的判定阈值为$0.04\,\mathrm{m}$。

图~\ref{fig:ch2-mpc-case-snapshots}展示了一个成功回合的可视化快照序列，呈现了上述FSM-MPC逻辑的闭环执行流程。该回合来自翻滚场景，共执行237步（$2.37\,\mathrm{s}$）后达到目标区域。接近阶段首先选择抓取面，并沿预抓取参考逐步靠近目标；末端到预抓取位姿的误差满足阈值后，状态机进入对齐阶段。对齐阶段中，MPC使用短时预测轨迹继续降低末端与抓取参考的相对误差；闭合判据连续满足3步后，控制器转入闭合。随后系统确认抓取稳定，之后进入搬运阶段，最终将立方体送入目标区域。

\begin{figure}[!htbp]
    \centering
    \captionsetup[subfigure]{font=footnotesize,skip=2pt}
    \begin{subfigure}[t]{0.48\textwidth}
        \includegraphics[width=\linewidth]{fig_ch2_mpc_case_00}
        \caption{初始}
    \end{subfigure}\hfill
    \begin{subfigure}[t]{0.48\textwidth}
        \includegraphics[width=\linewidth]{fig_ch2_mpc_case_01}
        \caption{选择抓取面}
    \end{subfigure}

    \vspace{0.45em}
    \begin{subfigure}[t]{0.48\textwidth}
        \includegraphics[width=\linewidth]{fig_ch2_mpc_case_02}
        \caption{接近20\%}
    \end{subfigure}\hfill
    \begin{subfigure}[t]{0.48\textwidth}
        \includegraphics[width=\linewidth]{fig_ch2_mpc_case_03}
        \caption{接近40\%}
    \end{subfigure}

    \vspace{0.45em}
    \begin{subfigure}[t]{0.48\textwidth}
        \includegraphics[width=\linewidth]{fig_ch2_mpc_case_04}
        \caption{接近60\%}
    \end{subfigure}\hfill
    \begin{subfigure}[t]{0.48\textwidth}
        \includegraphics[width=\linewidth]{fig_ch2_mpc_case_05}
        \caption{接近80\%}
    \end{subfigure}

    \vspace{0.45em}
    \begin{subfigure}[t]{0.48\textwidth}
        \includegraphics[width=\linewidth]{fig_ch2_mpc_case_06}
        \caption{进入对齐}
    \end{subfigure}\hfill
    \begin{subfigure}[t]{0.48\textwidth}
        \includegraphics[width=\linewidth]{fig_ch2_mpc_case_07}
        \caption{对齐中}
    \end{subfigure}
    \caption{FSM-MPC专家控制器在一个成功翻滚回合中的快照序列。快照按照实际闭环时间排序，展示从抓取面选择、接近预抓取位姿、短时MPC对齐、夹爪闭合、保持检测到最终搬运的完整状态机行为。}
    \label{fig:ch2-mpc-case-snapshots}
\end{figure}

\begin{figure}[!htbp]
    \ContinuedFloat
    \centering
    \captionsetup[subfigure]{font=footnotesize,skip=2pt}
    \begin{subfigure}[t]{0.48\textwidth}
        \includegraphics[width=\linewidth]{fig_ch2_mpc_case_08}
        \caption{闭合就绪}
    \end{subfigure}\hfill
    \begin{subfigure}[t]{0.48\textwidth}
        \includegraphics[width=\linewidth]{fig_ch2_mpc_case_09}
        \caption{进入闭合}
    \end{subfigure}

    \vspace{0.45em}
    \begin{subfigure}[t]{0.48\textwidth}
        \includegraphics[width=\linewidth]{fig_ch2_mpc_case_10}
        \caption{闭合中}
    \end{subfigure}\hfill
    \begin{subfigure}[t]{0.48\textwidth}
        \includegraphics[width=\linewidth]{fig_ch2_mpc_case_11}
        \caption{进入保持}
    \end{subfigure}

    \vspace{0.45em}
    \begin{subfigure}[t]{0.48\textwidth}
        \includegraphics[width=\linewidth]{fig_ch2_mpc_case_12}
        \caption{保持中}
    \end{subfigure}\hfill
    \begin{subfigure}[t]{0.48\textwidth}
        \includegraphics[width=\linewidth]{fig_ch2_mpc_case_13}
        \caption{进入搬运}
    \end{subfigure}

    \vspace{0.45em}
    \begin{subfigure}[t]{0.48\textwidth}
        \includegraphics[width=\linewidth]{fig_ch2_mpc_case_14}
        \caption{搬运中}
    \end{subfigure}\hfill
    \begin{subfigure}[t]{0.48\textwidth}
        \includegraphics[width=\linewidth]{fig_ch2_mpc_case_15}
        \caption{成功终止}
    \end{subfigure}
    \caption{FSM-MPC专家控制器在一个成功翻滚回合中的快照序列（续）。}
\end{figure}

\begin{table}[!htbp]
    \caption{FSM-MPC专家控制器的阶段逻辑。}
    \label{tab:expert-fsm-logic}
    \centering
    \small
    \renewcommand{\arraystretch}{1.15}
    \newcommand{\PhaseCell}[2]{\begingroup\renewcommand{\arraystretch}{0.82}\begin{tabular}[t]{@{}l@{}}#1\\[-0.45ex]{\footnotesize\textup{(#2)}}\end{tabular}\endgroup}
    \begin{tabularx}{\textwidth}{>{\raggedright\arraybackslash}p{0.13\textwidth}YY>{\centering\arraybackslash}p{0.12\textwidth}}
        \hline
        阶段 & 参考构造 & 切换条件 & 夹爪命令 \\
        \hline
        接近 & 跟踪预测抓取面的预抓取位姿；姿态先对齐夹爪接近轴，再对齐面内切向；接近过程中加入碰撞规避位移。 & 末端到预抓取参考的距离小于$0.04\,\mathrm{m}$时转入对齐；若候选抓取面不可达且无法切换则终止。 & 打开 \\
        对齐 & 在MPC时域内预测立方体短时运动，生成由预抓取位姿平滑过渡至抓取位姿的参考轨迹。 & 至少等待10步；位置误差小于$0.025+0.15\|v_c\|+0.02\|\omega_c\|$且姿态误差小于0.18，并连续满足3次后转入闭合；超过300步未满足则终止。 & 打开 \\
        闭合 & 保持抓取位姿参考，并在$0.30\,\mathrm{s}$内逐步闭合夹爪。 & 闭合时长满足要求且夹爪已闭合后转入保持。 & 闭合 \\
        保持 & 保持抓取位姿和闭合夹爪，检测物体是否随末端稳定保持。 & 末端与立方体中心距离连续20步小于$0.040\,\mathrm{m}$后转入搬运；长时间不稳定则终止。 & 闭合 \\
        搬运 & 将已抓取物体运送至目标区域；姿态跟踪权重降低，以优先保证平移运输稳定。 & 环境成功条件满足或回合超时。 & 闭合 \\
        \hline
    \end{tabularx}
\end{table}

对齐阶段的位置误差阈值随物体运动速度自适应调整：物体运动较快时，放宽位置容限以避免错过闭合时机；静止或低速时则收紧阈值，确保对齐精度。

MPC在每个控制步内求解一个短时域滚动优化问题，将任务空间参考转化为关节动作。设当前关节角为$q_0$，预测时域为$H$，优化变量为未来关节序列$q_{1:H}$，末端正运动学给出位置$p(q_h)$和姿态$R(q_h)$，任务空间参考为$(p^r_h,R^r_h)$。优化问题如下：
\begin{equation}
    \begin{aligned}
    \min_{q_{1:H}} \quad
    \sum_{h=1}^{H}\Big(
    & w_p\|p(q_h)-p^r_h\|_2^2
    + w_R\|e_R(R(q_h),R^r_h)\|_2^2 \\
    & + w_s\|q_h-q_{h-1}\|_2^2
    + w_v\left\|\frac{q_h-q_{h-1}}{\Delta t}\right\|_2^2
    + J_{\mathrm{lim}}(q_h)
    \Big),
    \end{aligned}
    \label{eq:expert-mpc-objective}
\end{equation}
其中$e_R(\cdot)$为姿态误差向量，$J_{\mathrm{lim}}$为关节限位惩罚项。优化同时约束关节位置和单步关节速度。

MPC预测时域设为$H=10$步，位置跟踪权重$w_p=300$显著高于姿态权重$w_R=50$以优先保证末端定位精度，平滑项权重$w_s=5$用于抑制关节抖动；最大迭代次数设为100。将上一时刻优化解在时域上向前滚动一步，作为当前优化问题的初值。在搬运阶段，控制器降低姿态误差和平滑项权重，以优先保证平移运输稳定。若MPC未收敛，或非搬运阶段跟踪误差超过阈值，则舍弃该关节序列，改用任务空间逆运动学作为替代控制方案。

该FSM-MPC控制器作为专家策略，为模仿学习提供示教数据。原始示教从翻滚场景采集，筛选出5000条成功专家轨迹用于训练。定轴自旋场景同样采集5000条示教，用作第四章行为克隆（BC）和扩散策略（DP）的验证对照（\S\ref{sec:main-clean}）。

数据集仅保留成功回合，并进一步筛选抓取相关阶段；动作标签取自专家参考轨迹。数据按每4步采样一帧的方式构建训练窗口，验证集占比5\%，数据划分采用固定伪随机种子以保证可复现性。主实验中，翻滚场景包含约19.4万个训练窗口和约1.0万个验证窗口。
